\chapter{Stima classica e limite di Cramér-Rao}

Nella stima bayesiana il parametro era una variabile aleatoria con una densità a priori nota. Spesso però non si ha
abbastanza informazione per assegnare una densità $f_\Theta(\theta)$: la massa di un elettrone, la distanza di una
stella, la probabilità di testa di una moneta mai vista non sono «estratte a caso» da nessuna parte. In questo capitolo
il parametro $\theta$ è una quantità \textbf{deterministica e non nota}: è l'inferenza \textbf{non bayesiana}, o
classica. Vedremo come costruire uno stimatore, come giudicarne la qualità e, soprattutto, che esiste un limite
inferiore all'errore che \emph{qualunque} stimatore può commettere: il limite di Cramér-Rao.

% ============================================================
\section{La verosimiglianza e lo stimatore a massima verosimiglianza}
% ============================================================

Le osservazioni $\vett{x} \in \mathcal{X}^n$ sono realizzazioni di un vettore aleatorio la cui PDF appartiene a una
famiglia $f_{\vett{X}}(\vett{x}; \theta)$, dove $\theta$ è deterministico e non noto e appartiene a un insieme $S$, lo
\term{spazio dei parametri}. Il punto e virgola, al posto della barra verticale del condizionamento, ricorda che $\theta$
non è una variabile aleatoria: è solo un parametro da cui la densità dipende.

\dfn{Verosimiglianza}{
  La \term{verosimiglianza} (\emph{likelihood}) del parametro $\theta$, date le osservazioni disponibili, è la funzione
  \[ L(\theta; \vett{x}) = f_{\vett{X}}(\vett{x}; \theta) , \]
  o, equivalentemente, la \term{funzione di verosimiglianza logaritmica} $\Lambda(\theta; \vett{x}) = \log
  f_{\vett{X}}(\vett{x}; \theta)$.
}

È la stessa formula della densità, guardata al contrario: la densità è una funzione dei dati $\vett{x}$ per $\theta$
fissato (e ha area 1); la verosimiglianza è una funzione del parametro $\theta$ per i dati fissati, quelli che si sono
osservati (e in generale non ha area 1). Dice quanto ciascun valore del parametro rende plausibili i dati osservati.

\dfn{Stimatore a massima verosimiglianza}{
  La \term{stima a massima verosimiglianza} di $\theta$ è
  \[ \hat\theta(\vett{x}) = \arg\max_{\theta \in S} \log f_{\vett{X}}(\vett{x}; \theta) , \]
  ed è una realizzazione dello \term{stimatore a massima verosimiglianza} (MLE, \emph{Maximum Likelihood Estimator})
  $\hat\Theta(\vett{X}) = \arg\max_{\theta \in S} \log f_{\vett{X}}(\vett{X}; \theta)$.
}

Si sceglie il valore del parametro che rende più probabili i dati osservati. Il logaritmo non sposta il punto di massimo
(Capitolo~21) e trasforma i prodotti delle densità di dati indipendenti in somme, molto più facili da derivare. Si noti
il legame con la stima bayesiana: con una densità a priori uniforme, la densità a posteriori è proporzionale alla
verosimiglianza, e la stima MAP coincide con quella a massima verosimiglianza.

\begin{figure}[H]
\centering
\begin{tikzpicture}
\begin{axis}[width=13cm, height=5.6cm, xmin=0, xmax=1, ymin=0, ymax=1.12, xtick={0,0.2,0.3,0.4,0.6,0.8,1},
  xticklabels={$0$,$0{,}2$,$0{,}3$,$0{,}4$,$0{,}6$,$0{,}8$,$1$}, ytick={0,0.5,1}, yticklabels={$0$,$0{,}5$,$1$},
  xlabel=$\beta$, ylabel={$L(\beta; \vett{x}) / L(\hat\beta; \vett{x})$}, axis lines=left,
  tick label style={font=\scriptsize}, label style={font=\small}, grid=major, grid style={mspigray!20},
  legend style={font=\scriptsize, at={(0.98,0.97)}, anchor=north east}]
  \addplot[mspiA, line width=1.1pt, domain=0.001:0.999, samples=200] {(x/0.3)^3*((1-x)/0.7)^7};
  \addplot[mspiB, line width=1.1pt, domain=0.001:0.999, samples=300] {(x/0.3)^30*((1-x)/0.7)^70};
  \draw[black, dashed] (axis cs:0.3,0) -- (axis cs:0.3,1.05);
  \legend{{$n = 10$, $w = 3$}, {$n = 100$, $w = 30$}}
\end{axis}
\end{tikzpicture}
\caption{La verosimiglianza della probabilità di testa $\beta$ di una moneta, divisa per il suo massimo, dopo $n$ lanci
con $w$ teste. In entrambi i casi il massimo è in $\hat\beta = \frac{w}{n} = 0{,}3$, ma con 100 lanci il picco è molto
più stretto: i dati escludono molti più valori del parametro.}
\end{figure}

% ============================================================
\section{Le qualità di uno stimatore}
% ============================================================

\dfn{Polarizzazione di uno stimatore}{
  Dato uno stimatore $\hat\Theta_n(\vett{X})$ del parametro non aleatorio $\theta$, calcolato da $n$ dati, si scrive
  \[ E\left[ \hat\Theta_n(\vett{X}) \right] = \theta + b_n(\theta) , \]
  dove $b_n(\theta)$ è la \term{polarizzazione} (\emph{bias}) dello stimatore. Lo stimatore è \term{non polarizzato}
  (o corretto) se $b_n(\theta) = 0$, e \term{asintoticamente non polarizzato} se $b_n(\theta) \to 0$ per $n \to \infty$.
}

Uno stimatore non polarizzato non sbaglia «in media»: se si potesse ripetere l'esperimento tante volte, la media delle
stime sarebbe il valore vero. Ma ogni singola stima può comunque essere lontana: serve anche misurare quanto le stime
sono sparpagliate. L'errore aleatorio dello stimatore si quantifica di solito con il suo valore quadratico medio,
l'\term{errore quadratico medio}
\[ e_n^2 = E\left[ \left( \hat\Theta_n(\vett{X}) - \theta \right)^2 \right] . \]

\teorema[msebias]{L'errore quadratico medio si divide in varianza e polarizzazione}{
  \[ \mathrm{MSE}\left[ \hat\Theta(\vett{X}) \right] = E\left[ \left( \hat\Theta(\vett{X}) - \theta \right)^2 \right] =
     \mathrm{VAR}\left[ \hat\Theta(\vett{X}) \right] + b_n^2(\theta) . \]
  In particolare, se l'errore quadratico medio tende a zero, tendono a zero sia la varianza sia la polarizzazione.
}

\begin{dimostrazione}
\passo{Si aggiunge e si toglie la media} $\hat\Theta - \theta = \left( \hat\Theta - E[\hat\Theta] \right) + \left(
E[\hat\Theta] - \theta \right)$, e la seconda parentesi è il numero $b_n(\theta)$.
\passo{Si sviluppa il quadrato}
\[ E\left[ \left( \hat\Theta - \theta \right)^2 \right] = E\left[ \left( \hat\Theta - E[\hat\Theta] \right)^2 \right] +
   2\, b_n(\theta)\, E\left[ \hat\Theta - E[\hat\Theta] \right] + b_n^2(\theta) . \]
Il termine misto ha media nulla, perché lo scarto dalla media ha media zero; il primo termine è la varianza.
\passo{La conseguenza} I due addendi $\mathrm{VAR}[\hat\Theta] \geq 0$ e $b_n^2(\theta) \geq 0$ non sono negativi: se
la loro somma tende a zero, tende a zero ciascuno dei due.
\end{dimostrazione}

\begin{figure}[H]
\centering
\begin{tikzpicture}[scale=0.62]
  \foreach \c/\t/\d in {0/{polarizzazione piccola\\ varianza piccola}/T1, 5.6/{polarizzazione piccola\\ varianza
    grande}/T2, 11.2/{polarizzazione grande\\ varianza piccola}/T3, 16.8/{polarizzazione grande\\ varianza grande}/T4} {
    \begin{scope}[xshift=\c cm]
      \fill[mspiB!10] (0,0) circle (2.2);
      \draw[mspigray] (0,0) circle (2.2) (0,0) circle (1.5) (0,0) circle (0.8);
      \fill[mspiB!70] (0,0) circle (0.22);
      \node[font=\scriptsize, align=center, anchor=north] at (0,-2.4) {\t};
    \end{scope}
  }
  \begin{scope}[xshift=0cm] \foreach \p in {(0.10,-0.03), (0.65,0.10), (0.04,0.08), (-0.10,-0.13), (0.10,-0.08), (0.15,-0.04), (-0.34,-0.25), (-0.16,-0.38), (-0.01,-0.05), (0.18,-0.16), (-0.04,0.11), (0.14,-0.19), (0.29,-0.12), (0.23,-0.06)} \fill[mspiA] \p circle (2.6pt); \end{scope}
  \begin{scope}[xshift=5.6cm] \foreach \p in {(-0.24,-0.72), (-0.07,1.38), (-0.37,0.89), (0.77,-0.22), (0.64,0.80), (-0.17,-0.31), (0.65,0.21), (1.23,0.25), (-0.66,-0.03), (-0.68,0.44), (0.41,-0.24), (-0.51,0.52), (-0.22,0.71), (-0.56,-0.16)} \fill[mspiA] \p circle (2.6pt); \end{scope}
  \begin{scope}[xshift=11.2cm] \foreach \p in {(0.65,0.25), (1.26,0.77), (0.63,0.60), (0.41,0.69), (0.96,0.39), (0.68,0.50), (0.84,0.38), (0.59,0.40), (0.67,0.60), (0.77,0.50), (0.44,0.38), (0.68,0.35), (0.73,0.39), (0.46,0.63)} \fill[mspiA] \p circle (2.6pt); \end{scope}
  \begin{scope}[xshift=16.8cm] \foreach \p in {(1.29,0.16), (1.12,-0.08), (0.03,0.34), (1.66,0.40), (-0.12,0.60), (0.43,0.05), (0.97,-0.54), (-0.27,-0.38), (1.39,0.52), (0.42,0.07), (1.15,1.19), (0.35,0.62), (0.58,0.67), (0.84,0.86)} \fill[mspiA] \p circle (2.6pt); \end{scope}
\end{tikzpicture}
\caption{Polarizzazione e varianza come tiri al bersaglio: il centro è il valore vero $\theta$, ogni punto blu una stima
ottenuta da un diverso campione. La polarizzazione è la distanza fra il centro dei colpi e il centro del bersaglio, la
varianza è quanto i colpi sono sparpagliati. Solo il primo stimatore ha un errore quadratico medio piccolo.}
\end{figure}

Per uno stimatore non polarizzato l'errore quadratico medio coincide con la varianza,
\[ \mathrm{VAR}\left[ \hat\Theta_n(\vett{X}) \right] = E\left[ \left( \hat\Theta_n(\vett{X}) - \theta \right)^2 \right]
   = E\left[ \hat\Theta_n^2(\vett{X}) \right] - \theta^2 , \]
e lo stimatore non polarizzato a minimo errore quadratico medio è quello che minimizza la varianza. Come nel caso
bayesiano (Capitolo~21), si parla poi di consistenza: uno stimatore è \term{consistente in senso debole} se
$\hat\Theta_n(\vett{X}) \to \theta$ in probabilità, \term{consistente in senso forte} se $\hat\Theta_n(\vett{X}) \to
\theta$ quasi certamente (con probabilità 1), \term{consistente in media quadratica} se $e_n^2 \to 0$.

% ============================================================
\section{Il limite di Cramér-Rao}
% ============================================================

Quanto può essere piccola la varianza di uno stimatore non polarizzato? Il risultato di questa sezione dice che non può
scendere sotto un certo valore, che dipende solo dalla famiglia di densità e non dallo stimatore. Sia $\vett{x} \in
\mathcal{X}^n$ la realizzazione di un vettore aleatorio $\vett{X} \sim f_{\vett{X}}(\vett{x}; \theta)$, con $\theta$
parametro non aleatorio. Supponiamo che la densità sia abbastanza regolare da poter derivare rispetto a $\theta$ sotto
il segno di integrale (Capitolo~21).

\begin{richiamo}{la derivata del logaritmo}
Per la regola della catena (Capitolo~11), la derivata di $\log g(\theta)$ è $\frac{g'(\theta)}{g(\theta)}$ (qui
$\log$ è il logaritmo naturale). Rovesciando la formula, $g'(\theta) = \left( \frac{\dd}{\dd\theta} \log g(\theta)
\right) g(\theta)$. Applicata a una densità che dipende da un parametro:
\[ \frac{\partial}{\partial\theta} f_{\vett{X}}(\vett{x}; \theta) = \left( \frac{\partial}{\partial\theta} \log
   f_{\vett{X}}(\vett{x}; \theta) \right) f_{\vett{X}}(\vett{x}; \theta) . \]
È il trucco che trasforma un integrale della derivata della densità in una \emph{media}, perché compare il fattore
$f_{\vett{X}}$.
\end{richiamo}

\teorema[score]{Le due identità della derivata logaritmica}{
  \[ E\left[ \frac{\partial}{\partial\theta} \log f_{\vett{X}}(\vett{X}; \theta) \right] = 0 , \qquad
     E\left[ \frac{\partial^2}{\partial\theta^2} \log f_{\vett{X}}(\vett{X}; \theta) \right] = -E\left[ \left(
     \frac{\partial}{\partial\theta} \log f_{\vett{X}}(\vett{X}; \theta) \right)^2 \right] = -\mathrm{VAR}\left[
     \frac{\partial}{\partial\theta} \log f_{\vett{X}}(\vett{X}; \theta) \right] . \]
}

\begin{dimostrazione}[derivando due volte l'area della densità]
\passo{Prima derivata} Si parte dall'identità $\int_{\R^n} f_{\vett{X}}(\vett{x}; \theta) \dd\vett{x} = 1$, che vale
per ogni $\theta$. Derivando rispetto a $\theta$ (il secondo membro è costante) e usando il richiamo:
\[ \int_{\R^n} \frac{\partial}{\partial\theta} f_{\vett{X}}(\vett{x}; \theta) \dd\vett{x} = \int_{\R^n} \left(
   \frac{\partial}{\partial\theta} \log f_{\vett{X}}(\vett{x}; \theta) \right) f_{\vett{X}}(\vett{x}; \theta)
   \dd\vett{x} = E\left[ \frac{\partial}{\partial\theta} \log f_{\vett{X}}(\vett{X}; \theta) \right] = 0 . \]
Sotto integrale $\vett{x}$ è una variabile muta di integrazione e si scrive con la lettera minuscola; nella media
statistica, che è la media di una funzione di un vettore aleatorio, si scrive con la maiuscola.
\passo{Seconda derivata} Si deriva di nuovo l'integrale del passo 1, con la regola del prodotto:
\[ \int_{\R^n} \left( \frac{\partial^2}{\partial\theta^2} \log f_{\vett{X}} \right) f_{\vett{X}} \dd\vett{x} +
   \int_{\R^n} \left( \frac{\partial}{\partial\theta} \log f_{\vett{X}} \right) \frac{\partial f_{\vett{X}}}{\partial
   \theta} \dd\vett{x} = 0 . \]
Nel secondo integrale si usa di nuovo il richiamo, $\frac{\partial f_{\vett{X}}}{\partial\theta} = \left(
\frac{\partial}{\partial\theta} \log f_{\vett{X}} \right) f_{\vett{X}}$: diventa la media del quadrato della derivata
logaritmica. Portandolo a destra si ottiene la seconda identità.
\passo{Quadrato medio e varianza} La derivata logaritmica ha media nulla (passo 1), quindi il suo valore quadratico medio
coincide con la sua varianza.
\end{dimostrazione}

\dfn{Informazione di Fisher}{
  L'\term{informazione di Fisher} che i dati contengono sul parametro $\theta$ è
  \[ I_n(\theta) = E\left[ \left( \frac{\partial}{\partial\theta} \log f_{\vett{X}}(\vett{X}; \theta) \right)^2 \right] =
     -E\left[ \frac{\partial^2}{\partial\theta^2} \log f_{\vett{X}}(\vett{X}; \theta) \right] . \]
}

La seconda forma ha un'interpretazione geometrica: la derivata seconda della log-verosimiglianza misura quanto è
\textbf{curva}, quanto è appuntito il suo picco. Un picco stretto, come quello con 100 lanci nella figura della prima
sezione, individua bene il parametro: tanta informazione. Un picco piatto lascia molti valori quasi ugualmente
plausibili: poca informazione.

\teorema[cramerrao]{Limite di Cramér-Rao}{
  Sia $\hat\Theta(\vett{X})$ uno stimatore del parametro non aleatorio $\theta$, con $E[\hat\Theta(\vett{X})] = \theta +
  b_n(\theta)$. Allora
  \[ \mathrm{VAR}\left[ \hat\Theta(\vett{X}) \right] \geq \frac{\abs*{1 + b_n'(\theta)}^2}{I_n(\theta)} . \]
  In particolare, se lo stimatore è non polarizzato,
  \[ \mathrm{MSE}\left[ \hat\Theta(\vett{X}) \right] = E\left[ \left( \hat\Theta(\vett{X}) - \theta \right)^2 \right] =
     \mathrm{VAR}\left[ \hat\Theta(\vett{X}) \right] \geq \frac{1}{I_n(\theta)} . \]
}

\begin{dimostrazione}[con la covarianza]
\passo{La media dello stimatore, derivata} Per ipotesi
$E[\hat\Theta(\vett{X})] = \int_{\R^n} \hat\Theta(\vett{x})\, f_{\vett{X}}(\vett{x}; \theta) \dd\vett{x} = \theta +
b_n(\theta)$. Derivando rispetto a $\theta$ (lo stimatore $\hat\Theta(\vett{x})$ dipende solo dai dati, non da
$\theta$) e usando il richiamo:
\[ \int_{\R^n} \hat\Theta(\vett{x}) \left( \frac{\partial}{\partial\theta} \log f_{\vett{X}}(\vett{x}; \theta) \right)
   f_{\vett{X}}(\vett{x}; \theta) \dd\vett{x} = E\left[ \hat\Theta(\vett{X})\, \frac{\partial}{\partial\theta} \log
   f_{\vett{X}}(\vett{X}; \theta) \right] = 1 + b_n'(\theta) . \]
\passo{È una covarianza} Ricordando che $COV(X, Y) = R_{XY} - \mu_X \mu_Y$ (Capitolo~9), la covarianza coincide con la
correlazione se almeno una delle due medie è nulla, come qui per la derivata logaritmica (Teorema~\ref{teo:score}):
\[ COV\left( \hat\Theta(\vett{X}),\ \frac{\partial}{\partial\theta} \log f_{\vett{X}}(\vett{X}; \theta) \right) =
   E\left[ \hat\Theta(\vett{X})\, \frac{\partial}{\partial\theta} \log f_{\vett{X}}(\vett{X}; \theta) \right] = 1 +
   b_n'(\theta) . \]
\passo{Il limite della covarianza} Per il Capitolo~9, $\abs{COV(X, Y)}^2 \leq \sigma_X^2 \sigma_Y^2$, quindi
\[ \abs*{1 + b_n'(\theta)}^2 \leq \mathrm{VAR}\left[ \hat\Theta(\vett{X}) \right]\, \mathrm{VAR}\left[
   \frac{\partial}{\partial\theta} \log f_{\vett{X}}(\vett{X}; \theta) \right] = \mathrm{VAR}\left[ \hat\Theta(\vett{X})
   \right] I_n(\theta) , \]
dove la varianza della derivata logaritmica è l'informazione di Fisher, per il Teorema~\ref{teo:score}.
\passo{Si divide} Dividendo per $I_n(\theta) > 0$ si ottiene il limite. Se lo stimatore è non polarizzato, $b_n = 0$,
la derivata $b_n'$ è nulla e la varianza coincide con l'errore quadratico medio.
\end{dimostrazione}

Mettendo insieme il limite con il Teorema~\ref{teo:msebias}, per uno stimatore qualsiasi
\[ \mathrm{MSE}\left[ \hat\Theta(\vett{X}) \right] = \mathrm{VAR}\left[ \hat\Theta(\vett{X}) \right] + b_n^2(\theta) \geq
   \frac{\abs*{1 + b_n'(\theta)}^2}{I_n(\theta)} + b_n^2(\theta) . \]
Il limite di Cramér-Rao (CRB, \emph{Cramér-Rao Bound}) è dunque un limite inferiore per l'errore quadratico medio di
qualunque stimatore non polarizzato: nessuno può fare meglio di $\frac{1}{I_n(\theta)}$.

\dfn{Stimatore efficiente}{
  Uno stimatore non polarizzato il cui errore quadratico medio è uguale al limite di Cramér-Rao, $\frac{1}{I_n(\theta)}$,
  si dice \term{efficiente}.
}

Le slide aggiungono due fatti, che non dimostriamo. Se per un dato problema di stima non bayesiana esiste uno stimatore
efficiente, questo coincide necessariamente con lo stimatore a massima verosimiglianza. E lo stimatore a massima
verosimiglianza ha un buon comportamento asintotico: sotto opportune condizioni di regolarità è
\begin{itemize}
  \item asintoticamente corretto (non polarizzato);
  \item consistente (converge in probabilità al valore vero);
  \item asintoticamente efficiente (la sua varianza si avvicina al limite di Cramér-Rao);
  \item asintoticamente normale (per $n$ grande ha densità gaussiana, come nel teorema centrale del limite).
\end{itemize}

% ============================================================
\section{Esempi}
% ============================================================

\begin{esempio}{La probabilità di testa}
Sia $\vett{x} \in \{0, 1\}^n$ la realizzazione di $n$ lanci di Bernoulli con parametro $\beta$ non noto, e sia
$w(\vett{x})$ il numero di «1». Come nel Capitolo~21, $p_{\vett{X}}(\vett{x}) = \beta^{w(\vett{x})} (1 - \beta)^{n -
w(\vett{x})}$; ora però $\beta$ non è aleatorio.

\textbf{La stima a massima verosimiglianza.} Si annulla la derivata della log-verosimiglianza, che è la stessa del
calcolo della stima MAP:
\[ \frac{\partial}{\partial\beta} \log p_{\vett{X}}(\vett{x}) = \frac{w(\vett{x})}{\beta} - \frac{n - w(\vett{x})}{1 -
   \beta} = 0 \quad\Longrightarrow\quad (1 - \beta)\, w(\vett{x}) = \beta\, (n - w(\vett{x})) \quad\Longrightarrow\quad
   \hat\beta = \frac{w(\vett{x})}{n} . \]
\textbf{Le prestazioni.} $w(\vett{X})$ è binomiale: $E[w(\vett{X})] = n\beta$ ed $E[w^2(\vett{X})] = n\beta(1 - \beta) +
n^2\beta^2$ (varianza più quadrato della media). Quindi
\[ E\left[ \hat\beta(\vett{X}) \right] = \frac{n\beta}{n} = \beta , \qquad E\left[ \hat\beta^2(\vett{X}) \right] =
   \frac{n\beta - n\beta^2 + n^2\beta^2}{n^2} = \frac{\beta(1 - \beta)}{n} + \beta^2 , \qquad \mathrm{VAR}\left[
   \hat\beta(\vett{X}) \right] = \frac{\beta(1 - \beta)}{n} . \]
(Nella slide 92 c'è scritto $\mathrm{VAR}[\hat\beta^2(\vett{X})]$: è la varianza di $\hat\beta$, non del suo quadrato.)
Lo stimatore è non polarizzato e consistente in media quadratica. \textbf{È efficiente?} La derivata seconda della
log-verosimiglianza è
\[ \frac{\partial^2}{\partial\beta^2} \log p_{\vett{X}}(\vett{x}) = -\frac{w(\vett{x})}{\beta^2} - \frac{n -
   w(\vett{x})}{(1 - \beta)^2} , \]
e quindi l'informazione di Fisher è
\[ I_n(\beta) = E\left[ \frac{w(\vett{X})}{\beta^2} + \frac{n - w(\vett{X})}{(1 - \beta)^2} \right] = \frac{n\beta}{\beta^2}
   + \frac{n - n\beta}{(1 - \beta)^2} = \frac{n}{\beta} + \frac{n}{1 - \beta} = n\, \frac{1 - \beta + \beta}{\beta(1 -
   \beta)} = \frac{n}{\beta(1 - \beta)} . \]
Il limite di Cramér-Rao è $\frac{1}{I_n(\beta)} = \frac{\beta(1 - \beta)}{n}$: la varianza dello stimatore a massima
verosimiglianza lo raggiunge, quindi lo stimatore è \textbf{efficiente}. Nessuno stimatore non polarizzato può stimare la
probabilità di testa meglio della frequenza delle teste.
\end{esempio}

\begin{esempio}{La media di dati gaussiani}
Siano $X_1, \ldots, X_n$ i.i.d. $\mathcal{N}(\mu, \sigma^2)$, con $\sigma^2$ noto e $\mu$ da stimare. La
log-verosimiglianza è
\[ \log f_{\vett{X}}(\vett{x}; \mu) = -\frac{n}{2} \log(2\pi\sigma^2) - \sum_{k=1}^{n} \frac{(x_k - \mu)^2}{2\sigma^2} ,
   \qquad \frac{\partial}{\partial\mu} \log f_{\vett{X}} = \sum_{k=1}^{n} \frac{x_k - \mu}{\sigma^2} , \qquad
   \frac{\partial^2}{\partial\mu^2} \log f_{\vett{X}} = -\frac{n}{\sigma^2} . \]
Annullando la derivata prima si trova $\hat\mu = \frac{1}{n}\sum_k x_k$: la stima a massima verosimiglianza è la media
campione. L'informazione di Fisher è $I_n(\mu) = \frac{n}{\sigma^2}$ e il limite di Cramér-Rao è $\frac{\sigma^2}{n}$, che
è proprio la varianza della media campione (Capitolo~18): anche qui lo stimatore a massima verosimiglianza è efficiente.
\end{esempio}

\gbox{In sintesi}{mspiA!80!black}{
\begin{itemize}
  \item Parametro \textbf{non aleatorio}: \textbf{verosimiglianza} $L(\theta; \vett{x}) = f_{\vett{X}}(\vett{x};
        \theta)$; stima a \textbf{massima verosimiglianza} $\hat\theta = \arg\max_\theta \log f_{\vett{X}}(\vett{x};
        \theta)$.
  \item \textbf{Polarizzazione} $b_n(\theta) = E[\hat\Theta] - \theta$; $\mathrm{MSE} = \mathrm{VAR} + b_n^2$; consistenza
        debole, forte, in media quadratica.
  \item Derivata logaritmica: media nulla; \textbf{informazione di Fisher} $I_n(\theta) = E[(\partial_\theta \log
        f)^2] = -E[\partial_\theta^2 \log f]$.
  \item \textbf{Limite di Cramér-Rao}: $\mathrm{VAR}[\hat\Theta] \geq \frac{\abs{1 + b_n'}^2}{I_n(\theta)}$; per uno
        stimatore non polarizzato $\mathrm{MSE} \geq \frac{1}{I_n(\theta)}$. Chi lo raggiunge è \textbf{efficiente}.
  \item Frequenza delle teste e media campione sono stime a massima verosimiglianza efficienti, con varianze
        $\frac{\beta(1-\beta)}{n}$ e $\frac{\sigma^2}{n}$.
\end{itemize}
}
